\section{Flexibility of auxiliary configurations}\label{sec:deformations}

The layer inequalities identify which constraints are tight and which
leave room for deformation. A \emph{contact} is a pair of distinct
points with inner product $1/2$. Here $X_{25}$ is the 197058-point
endpoint model of~\eqref{eq:d25}, and $X_T,X_M,X_R$ are the ERS layers
of Section~\ref{sec:codes}. The dimension-38 observation concerns
that layered-code construction; the 591900-point construction instead
uses the Leech lift of Section~\ref{sec:completion}.
The connection between contact
constraints and rigidity of spherical codes is developed
in~\cite{cjkt}; here it can be made quantitative directly from the
proved inner-product margins.

\subsection{Paired motion in dimension 25}
Write $s\in S_*/\sqrt{32}$ and
$x_s^{\pm}=(\sqrt3\,s/2,\pm1/2)$, with poles $p^{\pm}=(0,\pm1)$.
The contacts involving these new points are exactly the $496$ paths
\begin{equation}\label{eq:contact-paths}
 p^+\;--\;x_s^+\;--\;x_s^-\;--\;p^-.
\end{equation}
Their internal vertices are disjoint. They have no contacts with the
retained equator, because the relevant inner products are at most
$\sqrt3/4<1/2$. Between different directions the same-latitude bound
is $7/16$, and the opposite-latitude bound is $-1/16$. Thus the
displayed graph has $994$ vertices and $1488$ edges.

\begin{proposition}\label{prop:d25-motion}
Keep the equator and the two poles of $X_{25}$ fixed. For each
$s\in S_*/\sqrt{32}$ choose independently a unit vector $\widehat s$
with $\norm{\widehat s-s}<1/24$. Replacing $x_s^\pm$ by
$(\sqrt3\,\widehat s/2,\pm1/2)$ preserves the kissing condition and
the contacts~\eqref{eq:contact-paths}.
\end{proposition}
\begin{proof}
Let $\delta=\max_s\norm{\widehat s-s}<1/24$. For unit vectors,
\[
 \bigl|\ip{\widehat s}{\widehat t}-\ip{s}{t}\bigr|
 \le\norm{\widehat s-s}+\norm{\widehat t-t}\le2\delta.
\]
Consequently distinct same-latitude directions have inner product
less than $7/16+3\delta/2<1/2$, and distinct opposite-latitude
directions have inner product less than $-1/16+3\delta/2<0$.
For a fixed equatorial direction $z$,
\[
 \frac{\sqrt3}{2}\ip{z}{\widehat s}
 \le\frac{\sqrt3}{4}+\frac{\sqrt3}{2}\delta
 <\frac{13\sqrt3}{48}<\frac12.
\]
The last inequality follows from $3\cdot13^2<24^2$. A paired pair
still has inner product $1/2$, as do its contacts with the appropriate
poles. All other pole products are unchanged.
\end{proof}

With the equator, poles, and labels fixed, the common direction of
each pair is free to move in an open subset of $S^{23}$. The height is
forced, but the direction is not.
Thus saturation of the latitude constraints does not imply rigidity
of the lifted configuration.

\subsection{Independent motion of the dense layers}
In dimensions $m=n=38,39$, the dense sign codes have minimum distance
ten. Their mutual inner products are therefore at most $1-20/m<1/2$.
Their products with the middle and root layers are at most
$\sqrt{8/m}$ and $\sqrt{2/m}$, also strictly less than $1/2$.
Every dense point is an isolated vertex of the contact graph.

\begin{proposition}\label{prop:dense-motion}
For $m=n\in\{38,39\}$, keep $X_M$ and $X_R$ fixed and put
\[
 \eta_m=\frac{40-m}{4m}.
\]
Each $x\in X_T$ may independently be replaced by a unit vector
$\widehat x$ with $\norm{\widehat x-x}<\eta_m$. The resulting
configuration has the same cardinality and satisfies the kissing
condition. No new contact involves a dense point.
\end{proposition}
\begin{proof}
The change in a product of two unit vectors is at most the sum of
their displacement lengths. Hence distinct moved dense points have
inner product strictly less than
\[
 1-\frac{20}{m}+2\eta_m=\frac12.
\]
For a moved dense point and a fixed middle point, the bound is
$\sqrt{8/m}+\eta_m<1/2$. To verify the latter inequality exactly,
the positive number $1/2-\eta_m$ has square exceeding $8/m$ by
$153/5776$ at $m=38$ and by $937/24336$ at $m=39$.
The root-layer bound is smaller. The other pairs are unchanged.
\end{proof}

The respective displacement radii are $1/76$ and $1/156$.
The codes provide exact starting configurations, while their
neighbourhoods need not retain a binary description. The two motion
results exhibit distinct sources of flexibility: collective motion
of a contacting pair in dimension $25$, and independent motion of
strictly separated points in dimensions $38$ and $39$.
